% Part: lambda-calculus
% Chapter: church-rosser
% Section: definitions-and-properties

\documentclass[../../../include/open-logic-section]{subfiles}

\begin{document}

\olfileid{lam}{cr}{dap}

\olsection{Definisi lan Sipat}

Ing bab iki kita ngenalake
sipat Church--Rosser lan sawetara
akibat umume.

\begin{defn}[Sipat Church--Rosser, CR] Relasi $\xredone$ ing
  suku diarani nduweni \emph{sipat Church--Rosser}
  yen lan mung yen, saben
  $M \xredone P$ lan $M \xredone Q$,
  ana~$N$ kang nyukupi
  $P \xredone N$ lan $Q \xredone N$.
\end{defn}

Kalkulus lambda bisa dideleng minangka
model komputasi: suku kang wis ana
ing wujud normal iku ``nilai'',
dene relasi reduksi ing suku iku
``aturan pangétungan''.
Sipat Church--Rosser njamin yen
pangétungan bisa maju liwat luwih
saka siji dalan, saben asil kang
tekan wujud normal tetep padha.

Contone saka aljabar dhasar,
ana luwih saka siji cara
ngitung $4 \times (1+2) + 3$.
Ungkapan iki bisa direduksi dadi
$4 \times 3+3$ (yen luwih dhisik
$1+2$ direduksi dadi~$3$)
utawa dadi $4 \times 1+4
\times 2+3$ (yen luwih dhisik
$4 \times (1+2)$ diowahi
miturut sipat distributif).
Loro-lorone banjur bisa
direduksi dadi $12+3$.

Yen $\xredone$ dijupuk minangka
reduksi-$\beta$, akibat saka
sipat Church--Rosser yaiku
wujud normal sawijining suku,
yen ana, mesthi tunggal.
Upamane $M$ bisa direduksi
dadi $P$ lan $Q$, lan
loro-lorone wujud normal.
Miturut sipat Church--Rosser,
ana $N$ kang bisa digayuh
kanthi reduksi saka $P$
uga saka $Q$.
Amarga miturut pangira
$P$ lan $Q$ wujud normal,
reduksi saka $P$ lan $Q$
menyang $N$ mung bisa
reduksi sepele; tegese
$P$, $Q$, lan $N$ padha.
Iki mbenerake sebutan
\emph{wujud normal} sawijining suku.

Yen kalkulus lambda dideleng
minangka model komputasi,
wujud normal suku bisa dianggep
``asil pungkasan'' pangétungan
kang diwiwiti saka suku iku.
Akibat ing ndhuwur tegesé
asil pungkasan iku, yen ana,
mung siji, kaya pangétungan
$4 \times (1+2)+3$ kang
asile siji, yaiku~$15$.

\begin{thm} \ollabel{thm:str}
  Yen relasi $\xredone$ nduweni
  sipat Church--Rosser, lan
  $\xred$ yaiku relasi transitif
  paling cilik kang ngemot
  $\xredone$, mula $\xred$
  uga nduweni sipat Church--Rosser.
\end{thm}

\begin{proof}
  Upamane
  \begin{align*}
    M & \xredone P_1 \xredone \dots \xredone P_m \text{ lan}\\
    M & \xredone Q_1 \xredone \dots \xredone Q_n.
  \end{align*}
  Kita buktekake teorema iki
  kanthi mbangun kothak suku $N$
  kang dhuwure $m + 1$ lan
  ambane $n + 1$.
  $N_{i,j}$ nuduhake suku
  ing larik kaping-$i$ lan
  kolom kaping-$j$.
  
  Kita mbangun $N$ supaya
  $N_{i,j} \xredone N_{i+1,j}$
  lan $N_{i,j} \xredone N_{i,j+1}$.
  Ing pinggir kothak:
  \begin{align*}
    N_{0,0} &= M \\
    N_{i,0} &= P_i && \text{yen } 1 \le i \le m \\
    N_{0,j} &= Q_j && \text{yen } 1 \le j \le n \\
  \intertext{lan ing panggonan liyane:}
    N_{i,j} &= R
  \end{align*}
  kanthi $R$ suku kang nyukupi
  $N_{i-1,j} \xredone R$ lan
  $N_{i,j-1} \xredone R$.
  Ing saben sel njero, suku
  sisih ndhuwur lan kiwa iku
  asalé saka suku pojok
  ndhuwur-kiwa kanthi
  siji langkah.
  Sipat Church--Rosser kanggo
  $\xredone$ njamin yen suku
  kaya mangkono mesthi ana.

  Saiki $N_{m,0} \xredone \dots
  \xredone N_{m,n}$ lan
  $N_{0,n} \xredone \dots
  \xredone N_{m,n}$.
  Gatekna yen $N_{m,0}$
  yaiku $P_m$ lan $N_{0,n}$
  yaiku $Q_n$.
  Miturut definisi $\xred$,
  teorema iki kabukten.
\end{proof}

\end{document}
